\documentclass[11pt]{article}
\usepackage[margin=1in]{geometry}
\usepackage{amsmath,amssymb,amsthm,mathtools}
\usepackage{graphicx}
\usepackage{enumitem}
\usepackage{booktabs}
\usepackage[hidelinks]{hyperref}

\newtheorem{theorem}{Theorem}[section]
\newtheorem{proposition}[theorem]{Proposition}
\newtheorem{lemma}[theorem]{Lemma}
\newtheorem{corollary}[theorem]{Corollary}
\newtheorem{problem}[theorem]{Problem}
\newtheorem{remark}[theorem]{Remark}

\newcommand{\Baire}{\mathcal B_1}
\newcommand{\clp}{\operatorname{cls}_p}
\newcommand{\R}{\mathbb R}
\newcommand{\N}{\mathbb N}
\newcommand{\eps}{\varepsilon}
\newcommand{\spancl}{\overline{\operatorname{span}}}
\newcommand{\norm}[1]{\lVert #1\rVert}
\newcommand{\weakstar}{w^*}

\title{Corrected Audit Packet for Strongly Rosenthal Compacta}
\author{Run \texttt{fa\_banach\_001}: corrected partial}
\date{June 18, 2026}

\begin{document}
\maketitle

\begin{abstract}
This packet replaces the previous separable-case packet.  That proof is not
valid.  There are two independent failures: pointwise-convergent Baire-one
sequences may still be equivalent to the \(\ell_1\) unit-vector basis in the
supremum norm, and the proposed Banach-space quotient does not collapse the
fibers of the admissible cover.  The valid replacement is a topological
embedding of the target compactum into a dual space \(E^*\), where
\[
  E=\spancl^{\norm{\cdot}}\{\alpha(e(x)):x\in X\}\subset V^* .
\]
This is useful structure, but it is not enough to prove admissibility because
it does not identify \(E\) as the dual of a separable Rosenthal Banach space.
The original separable strongly-Rosenthal-versus-admissible problem therefore
remains open in this packet.
\end{abstract}

\section{Source question and corrected status}

The source is Glasner--Megrelishvili, \emph{Representations of dynamical
systems on Banach spaces not containing \(\ell_1\)}, arXiv:0803.2320
\cite{GM0803}.  It defines strongly Rosenthal compacta and admissible compacta
and asks whether the two classes coincide.

\begin{center}
\includegraphics[width=0.96\linewidth]{figures/open_problem_crop.png}
\end{center}

\noindent\textbf{Corrected run status.}  The previous active packet claimed to
prove that every separable strongly Rosenthal compactum is admissible.  That
claim should not be used.  The present packet records why the proof fails and
keeps the strongest valid replacement obtained from the audit.

\medskip
\begin{center}
\fbox{\parbox{0.92\linewidth}{
\textbf{Outcome.}  Corrected partial only.  No proof, counterexample, or
published resolution of the original question is recorded here.  The active
mathematical content is the embedding theorem of Section~\ref{sec:replacement}
and the two failure mechanisms in Sections~\ref{sec:l1failure} and
\ref{sec:quotientfailure}.}}
\end{center}

\section{Definitions and source tools}

Let \(X\) be compact metrizable.  A bounded family \(F\subset C(X)\) is called
Rosenthal when its pointwise closure satisfies
\[
   \clp(F)\subset \Baire(X).
\]
Equivalently, by the Rosenthal--Talagrand criterion used in the source, no
sequence in \(F\) is equivalent to the unit-vector basis of \(\ell_1\).

A compactum \(K\) is strongly Rosenthal if it embeds into \(\Baire(X)\) for
some compact metrizable \(X\).  It is admissible if it embeds into
\(\clp(F)\) for a bounded Rosenthal family \(F\subset C(X)\) on a compact
metrizable space.

The previous packet used two source results.

\begin{itemize}[leftmargin=2em]
\item Theorem 6.3 of \cite{GM0803}: for a Rosenthal family
      \(F\subset C(Y)\), the compactum \(L=\clp(F)\) admits a representation
      in \(V^{**}\) for a separable Rosenthal Banach space \(V\), with
      evaluation functionals \(\alpha(y)\in V^*\).
\item Theorem 6.16 of \cite{GM0803}: admissibility can be characterized via
      weak-star compact subsets of second duals of separable Rosenthal spaces.
\end{itemize}

\section{First failure: pointwise convergence does not prevent \(\ell_1\)}
\label{sec:l1failure}

The previous proof assumed that a pointwise-convergent sequence in
\(\Baire(X)\) cannot be equivalent, in the supremum norm, to the unit-vector
basis of \(\ell_1\).  This is false.

\begin{proposition}[A compact metrizable counterexample]
\label{prop:l1pointwise}
There is a compact metrizable space \(X\) and a pointwise-null sequence
\((f_n)\subset\Baire(X)\) such that \((f_n)\) is isometric to the unit-vector
basis of \(\ell_1\) in the supremum norm.
\end{proposition}

\begin{proof}
Let
\[
  X=\{\infty\}\cup\bigcup_{m=1}^\infty \{m\}\times\{-1,1\}^m
\]
be the one-point compactification of the countable discrete union of the
finite sign cubes.  This is compact metrizable.  Define
\[
  f_n(\infty)=0,\qquad
  f_n(m,\sigma)=
  \begin{cases}
  \sigma_n,& n\le m,\\
  0,& n>m .
  \end{cases}
\]
For each fixed point of \(X\), \(f_n\) is eventually zero, hence \(f_n\to0\)
pointwise.  Since \(X\) is countable compact metrizable, every real-valued
function on \(X\) is Baire-one; indeed every inverse image of an open interval
is countable and therefore \(F_\sigma\).

For scalars \(a_1,\ldots,a_N\), choose \(m=N\) and
\(\sigma_n=\operatorname{sgn}(a_n)\).  Then
\[
 \left|\sum_{n=1}^N a_n f_n(m,\sigma)\right|
 =\sum_{n=1}^N |a_n|.
\]
The reverse inequality is the triangle inequality.  Thus
\[
  \left\|\sum_{n=1}^N a_n f_n\right\|_\infty
  =\sum_{n=1}^N |a_n|.
\]
\end{proof}

\begin{corollary}
The graph-coordinate construction in the previous packet need not produce a
Rosenthal family.  Pointwise compactness of the target family in
\(\Baire(X)\) is not enough to rule out \(\ell_1\)-sequences in the
corresponding supremum-norm coordinate family.
\end{corollary}

\section{Second failure: the quotient does not collapse fibers}
\label{sec:quotientfailure}

The previous packet then attempted to repair the admissible cover by quotienting
the representing Banach space.  In the notation of Theorem 6.3, let
\[
  \nu_0:L\to V^{**},\qquad \alpha:Y\to V^*
\]
be the representation, and let \(e:X\to Y\) denote the graph embedding.  The
old proof set
\[
  N=\bigcap_{x\in X}\ker \alpha(e(x))\subset V
\]
and claimed that the quotient map \(q:V\to V/N\) makes
\(q^{**}\nu_0(h)\) depend only on the fiber of the restriction map
\(\pi:L\to K\).

There are two problems.

\begin{enumerate}[leftmargin=2em]
\item In the actual construction behind Theorem 6.3, the vectors represented
      by the graph coordinates generate the relevant predual space.  For the
      graph set, the above \(N\) is therefore typically just \(\{0\}\).  Then
      the quotient does not collapse any fibers.
\item The annihilator step uses the wrong closure.  If
      \(A=\{\alpha(e(x)):x\in X\}\subset V^*\), then
      \[
        \left(\bigcap_{\phi\in A}\ker\phi\right)^\perp
        =\overline{\operatorname{span}}^{\,\weakstar} A
      \]
      in \(V^*\), by the bipolar theorem for the dual pair \((V,V^*)\).
      But an element of \(V^{**}\) is norm-continuous on \(V^*\), not
      generally weak-star continuous.  Therefore vanishing on \(A\) gives
      vanishing on the norm-closed span of \(A\), not on its weak-star closed
      span.
\end{enumerate}

Thus the quotient proof cannot prove that \(K\) is weak-star compact in the
second dual of a separable Rosenthal quotient.

\section{The valid replacement}
\label{sec:replacement}

The norm-closure issue points to a genuine theorem.  It gives a topological
embedding into a dual Banach space, but not the second-dual representation
needed for admissibility.

\begin{theorem}[Norm-closed evaluation-span embedding]
\label{thm:embedding}
Let \(L\) be an admissible cover of a compactum \(K\), with continuous
surjection \(\pi:L\to K\), arising from a graph restriction construction.
Assume \(L\) is represented as in Theorem 6.3 of \cite{GM0803} by
\[
   \nu_0:L\to V^{**},\qquad \alpha:Y\to V^*,
\]
where \(V\) is separable Rosenthal and
\[
   h(y)=\langle \nu_0(h),\alpha(y)\rangle
   \qquad(h\in L,\ y\in Y).
\]
Put
\[
   E=\spancl^{\norm{\cdot}}\{\alpha(e(x)):x\in X\}\subset V^* .
\]
Then the formula
\[
   \sigma(\pi(h))=\nu_0(h)|_E
\]
defines a topological embedding
\[
   \sigma:K\hookrightarrow E^*
\]
when \(E^*\) carries its weak-star topology from the predual \(E\).
\end{theorem}

\begin{proof}
If \(\pi(h_1)=\pi(h_2)\), then
\[
 h_1(e(x))=h_2(e(x))\qquad(x\in X).
\]
Equivalently,
\[
 \langle \nu_0(h_1)-\nu_0(h_2),\alpha(e(x))\rangle=0
 \qquad(x\in X).
\]
Since \(\nu_0(h_1)-\nu_0(h_2)\in V^{**}\) is norm-continuous on \(V^*\), it
vanishes on the norm-closed linear span \(E\).  Hence
\(\nu_0(h_1)|_E=\nu_0(h_2)|_E\), so \(\sigma\) is well-defined.

Conversely, if the restrictions to \(E\) agree, then they agree on every
\(\alpha(e(x))\).  Therefore \(h_1(e(x))=h_2(e(x))\) for all \(x\), which is
precisely \(\pi(h_1)=\pi(h_2)\) in the graph-restriction quotient.  Thus
\(\sigma\) is injective.

Continuity follows because, for each \(u\in E\), the function
\[
  h\mapsto \langle \nu_0(h),u\rangle
\]
is continuous on \(L\) after first checking it on the generating
\(\alpha(e(x))\)'s and then passing to the norm closure uniformly on bounded
sets.  Therefore the factor map \(K\to E^*\) is weak-star continuous.  A
continuous injection from compact \(K\) into the Hausdorff space \(E^*\) is a
homeomorphism onto its image.
\end{proof}

\begin{remark}
Theorem~\ref{thm:embedding} is the strongest valid replacement found in this
audit.  It does not prove admissibility because it gives \(K\subset E^*\), not
\(K\subset W^{**}\) for a separable Rosenthal Banach space \(W\).  To apply
Theorem 6.16 of the source, one would need to identify \(E\) with \(W^*\), or
otherwise represent \(E^*\) as such a second dual in a way compatible with the
embedding.  No such identification follows from the construction.
\end{remark}

\section{Literature and remaining problem}

The audit supplied for this correction reports no published resolution of the
underlying open statement.  It also reports that later work continued to treat
the separable strongly Rosenthal case as unresolved: a later
Glasner--Megrelishvili survey still describes the separable strongly Rosenthal
case as unknown, Marciszewski--Pol formulate an equivalent envelope problem
and solve only the Aleksandrov--Urysohn special case, and specialist material
from 2021 continues to distinguish the relevant classes.

\begin{problem}
Does every separable strongly Rosenthal compactum admit an admissible
representation?
\end{problem}

\begin{problem}
Can the space \(E=\spancl^{\norm{\cdot}}\{\alpha(e(x)):x\in X\}\subset V^*\)
from Theorem~\ref{thm:embedding} be represented, under natural hypotheses, as
the dual of a separable Rosenthal Banach space?
\end{problem}

\begin{problem}
Can one build a strongly Rosenthal compactum for which the embedding of
Theorem~\ref{thm:embedding} cannot be lifted to a second-dual representation
over any separable Rosenthal space?
\end{problem}

\section{What this packet supersedes}

This packet supersedes the previous active packet in this folder.  The old
packet is preserved under
\[
 \texttt{history/invalid\_previous\_packet\_2026\_06\_18/}
\]
only as provenance.  Its claimed theorem that every separable strongly
Rosenthal compactum is admissible should not be cited as a result of this run.

\section{Human review checklist}

\begin{enumerate}[leftmargin=2em]
\item Check Proposition~\ref{prop:l1pointwise}: the compact space is metrizable,
      the functions are Baire-one, and the \(\ell_1\) estimate is exact.
\item Check that the old graph-coordinate argument really needs a
      no-\(\ell_1\) conclusion in supremum norm and therefore fails.
\item Check the annihilator closure distinction:
      weak-star closure is supplied by the bipolar theorem, while the fiber
      argument requires norm closure because elements of \(V^{**}\) are
      norm-continuous on \(V^*\).
\item Check that Theorem~\ref{thm:embedding} is well-defined and is a genuine
      topological embedding into \(E^*\).
\item Check that this embedding is insufficient for admissibility without an
      additional theorem identifying \(E\) as a Rosenthal dual predual.
\end{enumerate}

\begin{thebibliography}{9}

\bibitem{GM0803}
E. Glasner and M. Megrelishvili,
\newblock Representations of dynamical systems on Banach spaces not containing
\(\ell_1\),
\newblock arXiv:0803.2320.

\bibitem{Rosenthal}
H. P. Rosenthal,
\newblock A characterization of Banach spaces containing \(\ell_1\),
\newblock \emph{Proc. Nat. Acad. Sci. USA} \textbf{71} (1974), 2411--2413.

\end{thebibliography}

\end{document}
